\section{Limitations}\label{sec:limitations}

The main limitation of CliffordSTF is wall-clock cost: at matched parameter budget, a training step on rMD17 aspirin takes $3.39 \pm 0.01$\,s versus $1.44 \pm 0.02$\,s for the base Clifford model and $0.78 \pm 0.01$\,s for EquiformerV2 at $L{=}4$, roughly $2.4\times$ the base Clifford cost and slower than most spherical-harmonic baselines. The overhead stems from sequential dual-track message passing, the STF$_3$ generation step, and grade-sparse geometric-product dispatch in pure PyTorch without fused kernels; fused Cayley-table and STF kernels or sparse STF representations should narrow the gap, and the STF$_2$-only variant may suffice for long MD trajectories (Table~\ref{tab:ablation}). QM9 and Molecule3D additionally required gradient clipping to prevent early-training divergence from bilinear cross-track coupling, which a more principled cross-track initialisation or dedicated STF normalisation would likely address. Our controlled protocol fixes all models near $10^6$ parameters with a shared optimiser and learning rate, sacrificing per-model tuning for comparability; consequently MACE and ICTP---whose recommended recipes use learning rates two orders of magnitude above ours together with on-plateau scheduling---reach only 3000--6000\,meV/\AA{} force MAE on rMD17 versus tens of meV/\AA{} at recommended settings, and GotenNet, while matching our learning rate, additionally relies on $10{,}000$-step warmup and patience-150 early stopping that our protocol omits and similarly underconverges. Since published OC20/OC22 numbers use $10$--$100\times$ more parameters, our catalysis results characterise the sub-10M-parameter regime rather than the leaderboard frontier. Despite the order-of-magnitude relative improvement in rMD17 force-cosine similarity ($0.055 \rightarrow 0.551$), absolute directional fidelity remains below fully spherical-harmonic architectures (EquiformerV2 at $L\!\geq\!2$ attains $\sim\!0.96$); closing this remaining gap without spherical harmonics, Wigner-$D$ matrices, or CG tables is a natural extension. Finally, our directional-fidelity analysis is restricted to rMD17 small organics in vacuum: whether the $L \leq 1$ clustering at poor cosine alignment persists on periodic systems, larger molecules, or solvated condensed phases remains open, though CliffordSTF's best out-of-distribution and in-distribution S2EF energy MAE on OC22, together with the best $L\!\geq\!2$ IS2RE energy MAE in-distribution and out-of-distribution EwT in our study, suggest the phenomenon generalises to catalysis.

\section{Discussion}\label{sec:discussion}

This paper unifies two traditionally disjoint equivariant lineages. Geometric algebra treats the multivector as the primary substrate and derives rotational covariance from the geometric product, while irreducible Cartesian tensors treat symmetric traceless tensors as the primary substrate and derive covariance from algebraic symmetry. Our results suggest neither substrate alone is sufficient for atomistic force prediction at the per-edge bilinear level: the Clifford geometric product reaches $L \leq 1$ exactly and attains an aggregate rMD17 force cosine similarity of only $0.055$ under matched training; Cartesian symmetric traceless tensors cover $L \geq 2$ but forfeit the natural scalar and pseudoscalar structure of the Clifford algebra; the hybrid improves cosine similarity by an order of magnitude to $0.551$ while simultaneously improving force and energy mean absolute errors over the Clifford baseline. The broader implication is that representational choice in equivariant architectures should be governed by the angular content one actually wants to represent, not by loyalty to a particular algebraic formalism; Theorem~\ref{thm:cg_completeness} makes this concrete by showing that the geometric product and closed-form STF operations together span the same $\mathrm{SO}(3)$-irrep content as spherical harmonics through $L{=}3$. Our twelve-variant ablation isolates each track's contribution and confirms the two lineages are empirically complementary: a Clifford-only configuration (vanilla $L{=}1$ Clifford, no STF anywhere) reaches $f_{cos}{=}0.030$, an STF-only configuration (STF tracks at the output only, no STF inside the message-passing scaffold) reaches $f_{cos}{=}0.097$, and only the hybrid scaffold attains $f_{cos}{\in}[0.421, 0.506]$---an order-of-magnitude jump that neither track produces alone. The dominant driver is the $L \geq 2$ scaffold operating throughout message passing rather than any single architectural flag. Cross-track coupling acts as a training-stability lever rather than the main accuracy driver; this is consistent with the theoretical framing, since the scaffold is what unlocks representational capacity while the coupling governs whether the gradients flowing through it behave predictably. We see promising extensions along three axes: higher $L_{\max}$ through iterated STF constructions following \citet{shao2025orthonormal}, fused hardware kernels to close the wall-clock gap against spherical-harmonic baselines, and full $\mathrm{O}(3)$ equivariance through explicit parity tracking of polar and axial vector channels (Remark~\ref{rem:so3_vs_o3}).

\subsubsection*{Ethics Statement}
This work develops an architectural component for machine-learned interatomic potentials and involves no human subjects, personal data, or sensitive attributes; all benchmarks (rMD17, QM9, Molecule3D, OC20, OC22) are public and widely adopted. Dual-use considerations are those shared with computational chemistry generally and are low relative to existing quantum-chemistry software already in widespread use.

\subsubsection*{Broader Impact}
By closing the $L \geq 2$ gap of the geometric product without an e3nn-scale dependency stack, CliffordSTF lowers the barrier for practitioners who would benefit from equivariant interatomic potentials but lack specialised-library support. The broader "two lineages are complementary" principle should generalise to other $\mathrm{SO}(3)$-equivariant domains such as protein structure and fluid dynamics.

\subsubsection*{Reproducibility Statement}
We release the complete PyTorch implementation of CliffordSTF and every baseline, together with exact per-benchmark hyperparameters, at \texttt{https://anonymous.4open.science/r/Anon-CliffordSTF}; all training protocols are specified in \S\ref{sec:protocol} and the five ablation flags of Supplementary~\ref{supp:variants}, together with the structural \emph{STF only output} variant documented there, parameterise every variant used.